% ch-discretization.tex -- numerical method: the grid, the ghost layers and the matrix rows
\chapter{Discretization}
\label{ch:discretization}

This chapter describes the mesh Yuzen works on: where a quantity is allowed to
sit, how the boundary is padded, and how a point on the grid becomes a row of
the matrix. It is the vocabulary the rest of \cref{ch:timeintegration,%
ch:convergence,ch:linearsolver,ch:reproducibility} is written in, and it is the
first thing to read if you intend to add a physical term. The equations being
discretized are \cref{eq:height,eq:flux} of \cref{ch:governing}; the keys that
size the mesh are in \cref{sec:ref-domain:grid}.

\section{The staggered grid and the eight site types}
\label{sec:disc:stagger}

Yuzen uses a uniform rectangular grid. Given the grid lines
$x = x_{\min} + i_x\,\Delta x$, there are only two places along each axis at
which a quantity can be stored:

\begin{center}
\begin{tabular}{@{}l l@{}}
\code{0} & on the grid line, at $x = x_{\min} + i_x\,\Delta x$ \\
\code{1} & halfway between two grid lines, at $x = x_{\min} + i_x\,\Delta x + \Delta x/2$ \\
\end{tabular}
\end{center}

Three axes, two choices each: the list of possible positions is therefore
exactly three bits long and has exactly eight entries. Those three bits are the
type \code{Stagger<SX,SY,SZ>} (\file{yuzen/src/grid/Stagger.h}), and
\cref{tab:site-types} names them.

\begin{yzntable}{The eight site types and their three-bit encoding\label{tab:site-types}}{@{}l l L{0.46\textwidth}@{}}
\textsf{\bfseries Bits $(x,y,z)$} & \textsf{\bfseries Name} &
\textsf{\bfseries Position} \\ \midrule
0,0,0 & \code{Node}  & the grid point \\
1,0,0 & \code{Xedge} & the edge along $x$ \\
0,1,0 & \code{Yedge} & the edge along $y$ \\
0,0,1 & \code{Zedge} & the edge along $z$ \\
0,1,1 & \code{Xface} & the face with normal $x$ \\
1,0,1 & \code{Yface} & the face with normal $y$ \\
1,1,0 & \code{Zface} & the face with normal $z$ \\
1,1,1 & \code{Cell}  & the cell center \\
\end{yzntable}

The reason for encoding the positions this way rather than tabulating them is a
single observation: \textbf{moving half a cell along an axis flips that axis's
bit, and nothing else changes}. Every change of site type in the whole code is
that one operation, and the accompanying movement of the index follows from the
same bit in one line --- moving $+\tfrac12$ cell raises the index by one when
the bit was already \code{1}, and moving $-\tfrac12$ cell lowers it by one when
the bit was \code{0}. There is no hand-written table of offsets to keep in step
with anything.

In 1-D only the $x$ bit carries meaning, and in 2-D the $x$ and $y$ bits do.
The unknowns sit on \code{Xedge} in 1-D, on \code{Zface} in 2-D, and the
concentration of the \val{"CZ"} model (\cref{sec:conc:cz}) sits on \code{Cell}
--- in each case half a cell along every in-plane axis, which is what
\cref{ch:governing} calls the cell center.

A field is declared as \code{Field<GRID,S,T>}, with the site type \code{S} as a
template parameter rather than as runtime data. A stagger mismatch is therefore
a compile error, not a wrong answer: you cannot pass a cell-centered field where
an edge-centered one is expected. The overloads of
\code{C\_FiniteDifferences::GradX} and its relatives, one per site type, exist
for the same reason --- the derivative of a field knows which way it has moved
by half a cell, because the type says so.

\section{Ghost layers}
\label{sec:disc:ghost}

Outside the boundary the grid carries \textbf{two ghost layers on each side}
(\code{C\_Grid*::NumGhostX()}). Two, not one, because the thickness equation is
fourth order: evaluating it at a boundary point reaches two neighbors away.

\key{Grid.Nx} counts interior points only, so a grid declared with
\key{Nx = 300} holds $\mathrm{Nx} + 4 = 304$ points along $x$. Remember this
when you estimate memory or read a row count out of a debug print.

\code{CheckGhost(s)} sorts every point into one of three classes:

\begin{description}
  \item[\code{Inside}] within the computational domain. This is the only place
    a physical term adds a contribution.
  \item[\code{bXmin}, \code{bXmax}, \dots] the boundary itself, where the
    boundary condition fixes the value. In 2-D the corners join this set, and
    in 3-D the edges as well.
  \item[\code{Outside}] the outermost ghost layer.
\end{description}

\section{From a point to a matrix row}
\label{sec:disc:rows}

The number of rows is the \textbf{total number of points including the ghosts},
$\mathrm{imax()} - \mathrm{imin()} + 1$ --- not the number of interior points.
The map from a point to its row number is collected in one place per dimension,
\code{Index()} in \file{GridTraits.h}, and is given in \cref{tab:row-index}.

\begin{yzntable}{The point-to-row map, one per dimension\label{tab:row-index}}{@{}l l@{}}
\textsf{\bfseries Dimension} & \textsf{\bfseries Row} \\ \midrule
1-D & $\mathrm{row} = i_x$ \\
2-D & $\mathrm{row} = n_x\,i_y + i_x$ \\
3-D & $\mathrm{row} = n_x n_y\,i_z + n_x\,i_y + i_x$ \\
\end{yzntable}

Ghost rows are marked as boundary-condition rows with
\code{CRS\_Matrix::SetAsBC()}. A marked row is excluded from the time
integration --- \code{MakeMatrix\_IA()} and \code{MakeRHS()} skip it --- and the
physical terms consult \code{IsBC()} and do not add to it either. The
coefficients of a boundary-condition row are written once at initialization and
survive every subsequent zeroing of the matrix, so they never have to be
rebuilt.

The stencil is five points in 1-D, written \code{M m o p P}: the point itself
and two neighbors on each side. In 2-D it is thirteen points --- the
nine-point cross, plus the four diagonals \code{mm}, \code{pm}, \code{mp},
\code{pp}. \textbf{The diagonals are not optional.} They are there because the
biharmonic operator contains the cross term
$\partial^4/\partial x^2\partial y^2$, whose difference coefficient is of order
$1/(\Delta x^{2}\Delta y^{2})$ and whose stencil is precisely the four
diagonal neighbors. Drop them and you are no longer discretizing
$\nabla^{2}\nabla^{2}$.

The order in which the columns of a row are inserted is fixed by the ordering
of \code{T::FullIdx(k)} and its relatives. That ordering is not cosmetic: it
sets the order of the floating-point additions in the matrix--vector product,
and changing it breaks bit-level reproducibility. See \cref{ch:reproducibility}.

\section{The mobility at a face, and why its form is not free}
\label{sec:disc:positivity}

The film-height equation is a fourth-order \emph{degenerate} diffusion: the
mobility $D(h) = h^{3}$ vanishes where the film does. That degeneracy is what
makes a contact line possible at all, and it is also what makes a careless
discretization produce negative film thickness. A face between two cells needs
a single mobility built from the two thicknesses on either side, and the choice
of how to build it decides whether the scheme can go negative.

The well-mixed model uses the positivity preserving interpolation of Diez and
Kondic (2002), their Eq.~(23):
%
\begin{equation}
  D_{k}^{(x)} \;=\; \frac{2\,h_{k}^{2}\,h_{k+1}^{2}}{h_{k} + h_{k+1}},
  \qquad
  D_{k}^{(y)} \;=\; \frac{2\,h_{k}^{2}\,h_{k+n_x}^{2}}{h_{k} + h_{k+n_x}} .
  \label{eq:disc:pps}
\end{equation}
%
This is not an ad-hoc average. It is the general prescription
%
\begin{equation}
  D_{i} \;=\; \frac{h_{i+1} - h_{i}}{g_{i+1} - g_{i}},
  \qquad
  g(h) \;=\; \int \frac{\mathrm{d}h}{D(h)} ,
  \label{eq:disc:pps-general}
\end{equation}
%
evaluated for $D(h) = h^{s}$ at $s = 3$. The implementation is
\code{NonlinearTerm()} in \file{model/well\_mixed/LaplacePressureKernel.h}.

Why this form and not another: \cref{eq:disc:pps} factorizes as the product
$h_k h_{k+1}$ with the harmonic mean $2h_k h_{k+1}/(h_k + h_{k+1})$, so when one
side thins it vanishes as $O(h^{2})$, and it is exactly zero when one side is.
The face flux therefore switches itself off at the edge of a dry region. The
obvious alternative, the arithmetic mean $(h_k^{3} + h_{k+1}^{3})/2$, retains a
mobility of $h_k^{3}/2$ against a completely dry neighbor, draws liquid out of
the dry region, and drives $h$ negative.

\begin{warn}[Do not simplify this expression]
Rewriting \cref{eq:disc:pps} as an arithmetic mean --- or as anything else that
does not vanish with the thinner side --- removes the positivity guarantee. The
symptom is not a crash but negative film thickness near a contact line, which
then poisons the viscosity and the concentration equation downstream. The
property is pinned by
\file{yuzen/tests/test\_positivity\_mobility.cpp}, which checks
\cref{eq:disc:pps} against \cref{eq:disc:pps-general}, checks that the mobility
falls as $O(h^{2})$ and reaches zero, and checks that the arithmetic mean fails
both.
\end{warn}

\subsection{One face mobility for the terms of the potential}
\label{sec:disc:positivity-one}

The disjoining pressure and the vertical gravity take the same face mobility,
\cref{eq:disc:pps} to the bit (\code{Mobility::Face}; zero when either side is
dry or is the ghost $-h$ of a \val{"Pinning"} edge). The flux of a face is
then one mobility times the difference of the whole potential
$\varphi = \hat\sigma\nabla^2(H+S) + \Dpn\,\Pihat(H) - \Bo\,(H+S)$
(\cref{sec:output:potential}), so the dissipation of the scheme,
$\sum_f D_f(\Delta\varphi)^2/\Delta x$, is a sum of squares: the discrete free
energy never rises (\cref{sec:output:energy}), and a steady film is one with
$\varphi$ uniform, the minimum of the discrete energy. The vertical gravity's
term over a relief, $D_f\,\Delta s$, enters $A$ as $K(h_a + h_b)\,\Delta s$
with $K = D_f/(h_a + h_b)$, the same share for the two cells of the face. The
two terms took the arithmetic mean $(h_a^3 + h_b^3)/2$ until \file{Energy.dat}
showed the free energy rising near a steady film; that mean also left a
disjoining-pressure flux of $h^3/2$ against a dry neighbor.

\subsection{The face mobility with Navier slip}
\label{sec:disc:positivity-slip}

With \key{Physics.Slip} (\cref{sec:ref-physics:slip}) the mobility is
$D(h) = h^{3} + a h^{2}$, $a = 3\Lambda$, and \cref{eq:disc:pps-general}
applies with
%
\begin{equation}
  g(h) = -\frac{1}{a h} + \frac{1}{a^{2}}\ln\frac{h+a}{h}
       = -\frac{x - \ln(1+x)}{a^{2}}, \qquad x = \frac{a}{h} .
  \label{eq:disc:pps-slip}
\end{equation}
%
Both forms cancel somewhere, and \file{model/Mobility.h} switches to a series
there: $x - \ln(1+x)$ is summed as $\sum_{n=2}^{18}(-1)^n x^n/n$ below
$x = 0.1$ (remainder below $5\times10^{-19}$ relative), and where the two
heights differ by less than $10^{-3}$ of their mean $\bar h$ the face mobility
is $D(\bar h)/(1 + \delta^{2}K/24)$, $\delta = h_{k+1} - h_k$,
$K = D\,(1/D)''$, with a truncation error $O((\delta/\bar h)^4) < 10^{-12}$.
Against a reference free of cancellation the largest relative error is
$4\times10^{-13}$, and the value is continuous to $10^{-12}$ across each switch.
Without slip none of this is reached: \cref{eq:disc:pps} is
evaluated as written. The disjoining pressure and the vertical gravity take
the same face mobility with slip as without (\cref{sec:disc:positivity-one}).

The face mobility still vanishes on a dry side, but near it
$D \sim a h^{2}$, so it vanishes as $O(h)$ rather than $O(h^{2})$.

\begin{warn}[Slip weakens the positivity guarantee]
With slip the degeneracy that keeps the film positive is one order weaker.
Where the film is held up by a disjoining pressure, or does not thin below a
fraction of its height (about $0.1\,h_0$), this has no visible effect; a film
driven toward zero thickness with no disjoining pressure can go negative where
it would not without slip. Watch \code{HHmin} in \file{History.dat}.
\end{warn}

\subsection{The same mobility in the \texorpdfstring{\val{"CZ"}}{CZ} model}
\label{sec:disc:positivity-concz}

The \val{"CZ"} model resolves concentration, and hence viscosity, through the
thickness, so its in-plane velocity is not analytic: \code{CalVelocity()} in
\file{model/ConcZ/CZ2D/CZ\_HeightEquation\_2D.cpp} integrates $1/\mu$ and
$z/\mu$ across the film by Simpson's rule, and the concentration equation
advects with the result.

The film height itself, however, is advanced in the same conservative form as
the well-mixed model,
%
\begin{equation}
  \frac{\partial h}{\partial t} = -\frac{\partial Q}{\partial x} - E ,
  \label{eq:disc:cz-conservative}
\end{equation}
%
with the flux at a face carrying the same mobility \cref{eq:disc:pps}, divided
by three because the depth-integrated flux is $h^{3}/3\mu$ rather than $h^{3}$.
Integrating the \val{"CZ"} velocity profile through the thickness reproduces that
factor exactly in the uniform-viscosity limit. The implementation is
\code{FaceMobility()}, and the well-mixed counterpart of the same $1/3$ is the
\code{\_Magnitude} of its capillary term.

\begin{note}[The resolved viscosity at the face]
The flux through a face is the depth integral of the resolved velocity
profile,
$Q = h^3 b\int_0^1 (1-\zeta)^2/\mu\,d\zeta + h^2\tau\int_0^1 (1-\zeta)/\mu\,d\zeta$,
with $b$ the driving pressure gradient and body force and $\tau$ the
Marangoni stress. The mobility therefore carries the mean of $1/\mu(c(\zeta))$
weighted by $3(1-\zeta)^2$, the Marangoni flux the mean weighted by
$2(1-\zeta)$, and the surface tension is that of the surface concentration
(\code{FaceInvViscosity()}, \code{FaceSurfConc()}). The positivity-preserving
$h$ factors are unchanged, and a uniform viscosity gives $h^3/3\mu$ exactly.
A mobility that takes one viscosity and one surface tension per face at the
depth-mean concentration would move the film with a flux that is not the
integral of the profile which carries its solute, and with a stratified
concentration its error does not fall with refinement
(\cref{sec:verification:mms}).
\end{note}

Because both models are conservative and share the mobility, they agree in the
limit where the concentration is uniform through the thickness --- and that
agreement is a genuine check, since the two reach it by different routes. The
comparison is run by \file{tools/concz\_vs\_wm.py} and is part of the test
suite as \code{concz\_vs\_wm}: setting \key{Fluid.Viscosity.Exponent} to \val{0}
and switching off evaporation and solutocapillarity leaves the concentration
unable to affect $h$, and the peak film thickness of the two models then agrees
to within $10^{-6}$ relative, with the volume of each conserved exactly and
the same in the two \file{History.dat}.

\begin{note}[Two traps when running this comparison]
Leave the adaptive time step on and the two models choose different step
histories, so their output times drift apart; comparing the last row of each
\file{History.dat} then compares two different instants, which produces a
difference of a few percent that moves around non-monotonically under
refinement and looks exactly like a discretization error. Fix the step. Second,
the output interval must be well below \key{EndTime}, or the only shared
instant is the initial condition and the models appear to agree perfectly.
\end{note}

\section{The convection schemes}
\label{sec:disc:convection}

Three kinds of term are convective: the advection of the concentration, and
in the well-mixed film-height equation the in-plane body forces and the
Marangoni flux, which enter as a film height carried by a velocity rather than
as the gradient of a mobility. Each needs the value of a field at a face, and
takes it from the four cells $f_{LL}, f_L, f_R, f_{RR}$ about the face. For
the fixed schemes this value is
%
\begin{equation}
  f_{\text{face}} = \sum_j c_{1,j} f_j + \operatorname{sign}(u)\sum_j c_{2,j} f_j ,
  \label{eq:disc:face}
\end{equation}
%
with $c_1 = (-1, 7, 7, -1)/12$ and $c_2 = (-1, 3, -3, 1)/12$ for UTOPIA, and
$c_2$ three times larger for Kawamura--Kuwahara. Both are third order; the
leading error of the advection they give is $-|u|\Delta x^3/12\,\partial^4 f$
and $-|u|\Delta x^3/4\,\partial^4 f$. \val{"Central4"} has no $c_2$: its face
value is fourth order, the divergence only second, with no dissipation.
UTOPIA is the uniformly third-order polynomial interpolation of Leonard,
MacVean and Lock (1993), the Kawamura--Kuwahara scheme that of Kawamura,
Takami and Kuwahara (1986) written as a face value, and \val{"QUICK"} the
quadratic upstream interpolation of Leonard (1979); all are listed in Further
reading.

At the face of a \val{"Symmetry"} or \val{"Pinning"} edge the three
upwinded terms of the film-height equation carry no flux: their mobility there
is set to zero in $A$, in the Jacobian and in the velocity
(\code{UpwindKernelBase::WallFace}). Their flux is not a gradient, and a
mirrored ghost does not make the face value zero.

The limited schemes write the face value from the upwind side. With
$(UU, U, D)$ the cells $(LL, L, R)$ for a flow in $+x$ and $(RR, R, L)$ for one
in $-x$, and the two differences $a = f_U - f_{UU}$, $b = f_D - f_U$,
%
\begin{equation}
  f_{\text{face}} = f_U + \tfrac12\,\Phi(a, b), \qquad
  \Phi_{\text{Koren}} = \operatorname{sign}(b)\min\!\left(2|a|,\ \tfrac{2|b| + |a|}{3},\ 2|b|\right), \qquad
  \Phi_{\text{vL}} = \frac{2ab}{a + b},
  \label{eq:disc:limited}
\end{equation}
%
and $\Phi = 0$ unless $a$ and $b$ have the same sign: the limiters of Koren
(1993) and van Leer (1974). Nothing is divided by a gradient that may vanish. Koren's is UTOPIA wherever $0.4 \le a/b \le 4$ and
bounded elsewhere: the face value never leaves the interval between $f_U$ and
$f_D$, and at an extremum it falls back to $f_U$. That is the price of the
bound: a smooth extremum is clipped to first order locally, which costs the
accuracy of a smooth solution, and the maximum error converges at about 1.5
rather than 3 on a sine. The mirror image of a stencil gives the same value
to the last bit, so the checks of \cref{sec:verification:equivalence} hold for
every scheme.

In the concentration equation, which is linear in $c$ once its coefficients
are fixed, a limited face value is frozen at the latest iterate: in the row of
the face's upwind cell it is written $f_U + \chi/2\,(f_U - f_{UU})$ with
$\chi = \Phi/a$, elsewhere $f_U + \psi/2\,(f_D - f_U)$ with $\psi = \Phi/b$. The
frozen matrix is then an M-matrix, and the explicit half of Crank--Nicolson is
positive for a Courant number up to one, so the step keeps $c$ within its
bounds there. The step is not held to that; a step past it is reported in
\file{Warnings.txt}, for this equation and for the film height alike
(\cref{sec:ref-numerics:convection}). The well-mixed concentration is re-solved on every Newton pass,
so the frozen values converge with the film; the step is completed with the
limiter frozen at its own answer (\cref{sec:disc:wm-conservative}), and the
explicit half of each step is formed from the state the step starts from,
with the limiter at that state. In \val{"CZ"} each
stage uses the concentration it is evaluated at, which is the midpoint rule of
the limited operator.

\paragraph{$0 \le c \le 1$ without cutting or making solute.} The
concentration is a mass fraction, and each completed step is brought into
$0 \le c \le 1$. With the evaporation limited to the solvent there is
(\cref{sec:terms:evaporation}) only rounding passes 1 --- and the over- and
undershoot of an unbounded scheme at a sharp front: KK, UTOPIA, Central4 and
QUICK at any Courant number, Koren and van Leer above one. In the well-mixed
model what a real cell holds above 1, $h(c-1)$, goes to its face neighbors in
proportion to the room they have, $h(1-c)$; what they cannot take goes on in
a further pass from the cells it then sits in, up to four, and what is left
after them is spread over every cell with room, in proportion to it. A cell
below 0 is filled the same way the other way round: what it lacks, $-hc$, is
taken from its face neighbors in proportion to the solute $hc$ they hold, in
up to four passes, and then from every cell holding solute. The cells have
equal area, so the sum of $hc$ that \code{ResinVolErr} balances
(\cref{sec:output:history}) is kept to rounding, and the ghosts are formed
again from the cells as moved. Each pass takes the room of every cell before
it moves anything, so the result does not depend on the order of the cells,
and where no cell leaves $[0, 1]$ nothing is touched. Only a film holding more
solute than liquid in all, or less than none, would leave a remainder; it
would be cut, or made up, and reported in \file{Warnings.txt}.
\file{JobInfo.txt} gives the solute moved over the run
(\cref{sec:output:jobinfo}). Left below 0, the undershoot of a front from
$c = 0.7$ to $0.005$ reached $-0.04$ to $-0.09$ and, past the divergence test
of $-0.1$, stopped the run with QUICK and Central4, and with UTOPIA at
Courant 1.5 (found; corrected). Within the Newton passes the working field from
which the viscosity, the Marangoni stress and the evaporation are taken is
capped at 1, but the limiter of Koren and van Leer is frozen at the iterate as
the equation solved it: frozen at the capped values, the two rows that share a
face would take different face values from it, and the advection would lose
its flux form there. In \val{"CZ"} the excess of a cell goes to the cell
below it in its column, whose cells have the same volume; a column whose mean
passes 1, which the solvent limit prevents, keeps $c$ above 1 with a warning.
A cell below 0 is filled from the cells of its column that hold solute, in
proportion to it; a column whose mean is below 0, which is what the
undershoot of the advection along the film makes, takes what it lacks from
the columns next to it in proportion to their solute (each scaled as a whole,
so no cell of theirs goes below 0), in up to four passes and then from every
column holding solute, as the well-mixed cells do
(\code{CZ\_HeightEquation\_2D::RedistributeDeficit}). Left below 0, a front
of \val{"CZ"} from 0.7 to 0.005 with UTOPIA at Courant 0.4 passed the
divergence test and stopped the run (found; corrected).

In the film-height equation the matrix $A$ holds the face weights frozen at
the current film, so that $A(h)h$ is the flux exactly, and the Jacobian adds
the rest of the derivative of the limited face height $h_f$: with
$g = \partial h_f/\partial h$, $w(3g - w_t)h_f^2/\mu$ for the body forces and
$(2g - w_t)h_f\,\Delta c/\mu$ for the Marangoni flux. Newton keeps its
quadratic convergence (\cref{sec:verification:schemes}); at the kinks of
Koren's $\Phi$ the derivative is one-sided. The Marangoni flux is upwinded by
the direction of the flux itself, the sign of $\sigma_c\,\Delta c$, not of
$\Delta c$ alone, which would pick the downwind side with a negative
\key{Fluid.SurfaceTension.Coef}.

\subsection{The well-mixed concentration in conservative form}
\label{sec:disc:wm-conservative}

The well-mixed concentration equation is solved as the finite-volume balance
of the solute $hc$ of each cell over the step,
%
\begin{equation}\label{eq:disc:wm-conservative}
  \begin{aligned}
    h^{1}c^{1} - h^{0}c^{0}
      = {}& -\Delta t \sum_f s_f \bigl[\theta\,G_f^{1} + (1-\theta)\,G_f^{0}\bigr]
        + \Delta t\,S\,c_S \\
      &{} - \Delta t\,\bigl[\theta\,r^{1}c^{1} + (1-\theta)\,r^{0}c^{0}\bigr],
    \qquad
    G_f = F_f\,c_f - D(c_f)\,h_f\,\frac{c_{\mathrm{hi}} - c_{\mathrm{lo}}}{\Delta x},
  \end{aligned}
\end{equation}
%
$s_f = \pm 1/\Delta x$ the outward sign of face $f$ over the spacing across
it, $\theta = 1/2$. $F_f$ is the volume flux the film-height equation puts
through that face: the sum, over the enabled terms, of what each holds in $A$
for it --- the capillary, vertical-gravity and disjoining fluxes, and the
in-plane body forces, the centrifugal force and the Marangoni stress with
their upwinded face heights (\code{FaceFlux} of each kernel). $c_f$ is the
face value of \key{Numerics.Convection.Concentration} (\cref{eq:disc:face}),
upwinded by the sign of $F_f$; $D(c_f)\,h_f$ takes the means of the two cells.
$S$ is the liquid the nozzle supplies, at the middle of the step as the
film-height equation takes it, and $c_S$ the concentration it arrives at:
$c_S = \theta c^{1} + (1-\theta)c^{0}$, the concentration it lands on, with
\key{Physics.Source.ConcType = "local"} (the default), and
$c_S = \key{Physics.Source.Conc}$ with \val{"supplied"}, which mixes the new
liquid through the thickness of the cell's column: on a film with no flow,
$c^{1} = (h^{0}c^{0} + \Delta t\,S\,c_S)/(h^{0} + \Delta t\,S)$.
$r$ is the sideways outflow of the 1-D centrifugal term, which
leaves at the concentration of its cell. Evaporation takes solvent only and
does not appear: it is in $h^{1}$. This is \cref{eq:conc-dim} itself; the form
of \cref{eq:conc-adv} for $c$ agrees with it in the continuum, but discretely
only the balance keeps the solute. Solved for $c$ in the form of
\cref{eq:conc-adv}, a front from 0.9 to 0.3 on a closed
periodic film, carried by the flow and drying by Raoult's law, lost
$1\times10^{-5}$ to $3\times10^{-4}$ of its solute, depending on the scheme
and the grid (50 to 200 cells), over 0.3 of a crossing of the domain (found;
corrected: now $4\times10^{-16}$ to $3\times10^{-15}$ in 1-D and
$2\times10^{-14}$ in 2-D, \code{conc\_drying} in
\cref{sec:verification:analytic}).

The equation is solved for the change $d = c^{1} - c^{0}$ of the step. Taking
$c^{0}$ times the film's own step,
$h^{1} - h^{0} = -\Delta t \sum_f s_f[\theta F_f^{1} + (1-\theta)F_f^{0}]
- \Delta t\,[\theta r^{1} + (1-\theta)r^{0}] - \Delta t\,(E - S)$, from
\cref{eq:disc:wm-conservative} leaves
%
\begin{equation}\label{eq:disc:wm-change}
  \begin{aligned}
    &h^{1} d + \theta\Delta t \sum_f s_f\bigl[F_f^{1} d_f - K_f^{1}(d_{\mathrm{hi}} - d_{\mathrm{lo}})\bigr]
       + \theta\Delta t\,(r^{1} - S)\,d \\
    &\quad= -\Delta t \sum_f s_f \bigl[\theta F_f^{1} + (1-\theta)F_f^{0}\bigr](c^{0}_f - c^{0})
       + \Delta t\,E\,c^{0} \\
    &\qquad\quad{}+ \Delta t \sum_f s_f \bigl[\theta K_f^{1} + (1-\theta)K_f^{0}\bigr](c^{0}_{\mathrm{hi}} - c^{0}_{\mathrm{lo}}),
  \end{aligned}
\end{equation}
%
$K_f = D(c_f)h_f/\Delta x$, every face value $c^{0}_f - c^{0}$ formed as the
sum of the face weights times $c^{0}_k - c^{0}$. The weights sum to one, so a
uniform $c^{0}$ gives a right-hand side of exactly zero without evaporation,
and $d = 0$; the evaporation, $\Delta t\,E\,c^{0}$, concentrates the solute
as $h^{1}c^{1} = h^{0}c^{0}$ where nothing else acts. With
\key{EnableConvection} off the terms in $c_f - c$ are dropped, and the film
still thins or thickens under the solute ($\theta\Delta t \sum_f s_f
F_f^{1} d$); the solute is then not a balance, since the flow moves liquid
without the solute in it, and the step ends with $c^{0} + d$.

The film $h^{1}$ comes from the Newton iteration of the same step. On every
pass the concentration is solved with the film of that pass and its fluxes
(\cref{sec:time:coupling}); these give that film only once the iteration has
converged. The step is then completed (\code{AfterNKLoop} of
\file{HeightEquationBase.h}). The film is the one the iteration converged
to; the fluxes, the sideways outflow and the supply are taken at it, with the
concentration of the last pass, and the evaporation as that pass took it.
What their divergence still misses of the step the film took, cell by cell,
is the residual of the step at that film; less its mean over the cells, which
is the volume error of the step and stays as that, it is added to the fluxes
of the inner faces as running sums --- along $y$ the sum of each row of cells
spread evenly over the row, then along $x$ within each row what is left of
it --- so that the face fluxes are exactly those of the step the film took
(\code{ConsistentFluxes}). No edge face and no wall is touched, and nothing
is solved. The film the fluxes would give instead, $h^{0}$ less their
divergence, differs from the converged one by that residual, and on a film far
thicker than its cells the residual is at the rounding of fluxes of enormous
size ($\epsmach\,\Delta t/\Delta t_{\mathrm{explicit}}$,
\cref{sec:convergence:floor}): taken as the film, it goes negative in
\file{nozzle\_1d\_rotating}. The concentration is then
solved once more with these fluxes and this film, and the solute of each cell
is formed from \cref{eq:disc:wm-conservative} with the solute flux of each
face formed once, from the four cells about it whichever side asks
(\code{Reconstruct} of \file{ConcEquationBase.h}): the sum of $hc$ over the
cells changes by the fluxes through the edges and the sources alone, to
rounding, whatever the tolerance of the linear solve. Those edge fluxes, the
volume flux of each edge face (which \code{ConsistentFluxes} leaves alone)
and the solute flux \code{Reconstruct} formed there, are what
\file{History.dat} counts as crossing a \val{"FreshInlet"} or an
\val{"Outlet"} (\cref{sec:output:history}). A limited scheme is
solved again with its limiter frozen at its own answer until that settles
($10^{-12}$, at most eight times), so that the face values of the step are
the limited face values of the answer. Then $c$ is brought into $[0, 1]$
(\cref{sec:disc:convection}). A pass, or the end of a step, whose film has a
cell without liquid ($h \le 0$, at a step far too long for the film) leaves
the concentration as the pass before had it: the balance of $hc$ is not one
there, and the step is left to fail on the film (found: the
solve failed first and the run stopped naming the concentration solver;
corrected). At such a step the Newton correction of the film stagnates far
from convergence, and the step is cut, or at a fixed step the run stops
naming the stagnation (\code{wm\_diverges} in \file{tools/input\_guards.py},
\cref{sec:convergence:stagnation}); a film that is below zero all the same,
as one read from a file, stops the run as diverged (\code{wm\_film\_below}).

A uniform concentration under a strong in-plane flow, a nozzle, the vertical
gravity, the capillary pressure and, in 1-D, the spin stays uniform to
$6\times10^{-16}$ (\code{conc\_uniform}) --- in general to the rounding of
the solute fluxes relative to the film, $\epsmach\,\Delta t\sum_f|F_f|/h$,
which the solute formed from them once more carries (up to $3\times10^{-12}$
in \file{tools/proptest.py}, on a thin film next to a \val{"Pinning"} edge); with every term of the model the
manufactured solutions converge at second order in $\Delta x$ and $\Delta t$
(\cref{sec:verification:mms}), and the oblique wave of
\cref{sec:verification:schemes} at 1.91 and 1.97.

\subsection{The concentration of \texorpdfstring{\val{"CZ"}}{CZ} in conservative form}
\label{sec:disc:cz-conservative}

The \val{"CZ"} concentration equation is a finite-volume balance of $hc$ on
the $\zeta$ grid. The film-height equation stores the flux through each of
its faces and writes the end of the step as their divergence; the
concentration equation splits each of these into layer fluxes, the cell mean
of the resolved velocity profile adjusted so that the layers of a column add
up to the film-height flux face by face. The flux through the $\zeta$ faces
then follows from the discrete conservation of the grid, built up from the
substrate, with the evaporating solvent leaving through the surface face.
The ghost cells above the surface and below the substrate are set so that
the total flux there --- advection and the whole diffusive flux, its part
along the slope included --- is zero. The step is solved for the change of
$c$, so a uniform $c$ gives exactly no change, and the small solute residual
of the linear solver is put back as a uniform shift of the change, sized for
the solvent the shift itself sends through the surface term and applied to
the ghosts as well (each side of a periodic seam forms the flux there from
its own), so that the solute of the step balances to rounding. The limited schemes keep
their frozen upwind and downwind forms, which keep $c$ bounded, and are
re-solved until the limiter has settled at the midpoint of the step.

The scheme of \key{Numerics.Convection.Concentration} is used on every face,
$x$, $y$ and $\zeta$ alike. With the central schemes nothing is dissipated in
any direction, so the solution has what the model has at the shortest waves
the grid resolves: a film stratified with $\sigma_c\,\partial c/\partial\zeta
> 0$ and a weak diffusivity is unstable there, and a central scheme lets
those waves grow where an upwind-biased one damps them on a coarse enough
grid (\cref{sec:verification:schemes}).
